\documentclass{chsmc}
\chsmcsetup{
	year = 2025,
	part = 1,
	type = B,
	show-answers = false,
}
\begin{document}

\section{填空题：本大题共 8 小题，每小题 8 分，满分 64 分.}

\begin{enumerate}
	\item 设集合 $A = \{a^2 \mid a \in \mathbf{Z},\, a^2 \geqslant 2025\}$，
		$B = \{2^b \mid b \in \mathbf{Z},\, 9 \leqslant b \leqslant 14\}$，
		则 $A \cap B$ 的元素个数为 \blankline
	\item 若 $\sin 20^\circ \sin 25^\circ + \sin \alpha = \cos 20^\circ \cos 25^\circ$，
		则 $\cos 2\alpha$ 的值为 \blankline
	\item 若 $\log_3 x$, $\log_9(3x)$, $\log_{27}(9x)$ 成等比数列，
		则正数 $x$ 的值为 \blankline
	\item 设复数 $z$ 满足 $z+\ii$ 与 $z^2+2\ii$ 均为实数（$\ii$ 为虚数单位），
		则 $\bigl|z^3+3\ii\bigr|$ 的值为 \blankline
	\item 若 $x$, $y > 0$ 且 $x+y=1$，则 $\dfrac{y+1}{x} + \dfrac{x+2}{y}$ 的最小值为 \blankline
	\item 设 $F_1$, $F_2$ 分别为椭圆 $\Gamma\colon \dfrac{x^2}{a^2} + \dfrac{y^2}{b^2} = 1$
		（$a>b>0$）的左、右焦点，若 $\Gamma$ 上存在一点 $P$，使得直线 $F_1P$, $F_2P$
		的斜率分别为 $\dfrac{1}{2}$, $2$，则 $\Gamma$ 的离心率为 \blankline
	\item 平面中的 $3$ 个单位向量 $\vec{a}$, $\vec{b}$, $\vec{c}$ 满足
		$\vec{a}\cdot\vec{b} \geqslant \dfrac{1}{2}$，$\vec{b}\cdot\vec{c} \leqslant -\dfrac{1}{2}$，
		则 $\bigl|\vec{a}+\vec{b}+\vec{c}\bigr|$ 的最大值与最小值之和为 \blankline
	\item 从 $20$ 个数 $1$, $2$, $3$, $\cdots$, $20$ 中选出 $4$ 个不同的数（不计顺序），
		使它们的乘积为 $2025$ 的倍数，则不同选法的数目为 \blankline
\end{enumerate}

\section{解答题：本大题共 3 小题，满分 56 分. 解答应写出文字说明、证明过程或演算步骤.}

\begin{enumerate}[resume]
	\item (本题满分 16 分) 设 $f(x)$ 是定义域为 $\mathbf{R}$ 的函数，
		$g(x) = (x-1)f(x)$，$h(x) = f(x) + x$.
		若 $g(x)$ 为奇函数，$h(x)$ 为偶函数，求
		$\dfrac{f(1)f(3)\cdots f(99)}{f(2)f(4)\cdots f(100)}$ 的值.
	\item (本题满分 20 分) 在平面直角坐标系中，一条过点 $(0,1)$ 的直线 $l$ 经过点集
		$\Gamma = \{(x,y)\mid y^2 = x + 3|x|\}$ 中的四个点 $(x_i,y_i)$（$i = 1,2,3,4$）.
		求 $\dfrac{1}{x_1} + \dfrac{1}{x_2} + \dfrac{1}{x_3} + \dfrac{1}{x_4}$ 的取值范围.
	\item (本题满分 20 分) 对整数 $n \geqslant 3$，在一个棱长均为 $1$ 的正 $n$ 棱柱的所有 $3n$
		条棱中，随机选取两条不同的棱 $l_1$, $l_2$，将事件“$l_1$ 所在直线与 $l_2$ 所在直线平行”
		发生的概率记为 $P_n$.
		是否存在两个不同的正整数 $k$, $l$（$k,l \geqslant 3$）满足 $P_k = P_l$？证明你的结论.
\end{enumerate}

\end{document}
